\section{Restriction, extension and component boundaries}
\label{sec:recognition}
An incidence system in this paper is a nonempty finite simple graph $X$, together
with a surjective proper coloring $t:V(X)\to I$. Thus adjacent vertices have
different colors. A \emph{correlation} is a graph automorphism $f$ that permutes
the color classes: there is a permutation $\sigma_f$ of $I$ such that
$t(f(v))=\sigma_f(t(v))$ for all $v$. Surjectivity makes $\sigma_f$ unique.
Write $C(X)$ for the correlation group and $K(X)$ for its subgroup preserving
each color. The map $f\mapsto\sigma_f$ is a homomorphism, so $K(X)$ is normal
and $C(X)/K(X)$ embeds in $\Sym(I)$. In general $C(X)$ need not be the full
uncolored graph automorphism group.

A \emph{chamber} is a clique meeting every color class. An incidence system is
an \emph{incidence geometry} when each clique, including the empty clique, is
contained in a chamber. No transitivity, thinness or residual connectedness is
part of this definition. These are the conventions of
D\'iaz Ramos--Molinier--Viruel~\cite{DMV}, Definitions 2.1--2.3.

The following elementary fact is the common group-theoretic argument. It is
stated for arbitrary permutation groups; finiteness enters the graph
constructions later.

\begin{theorem}[Restriction and extension]\label{UR.Restriction}
Let $\iota:V\hookrightarrow W$ be an injection, let
$A\leq\Sym(V)$ and $B\leq\Sym(W)$, and let $A_0\leq A$ and $B_0\leq B$.
Suppose that:
\begin{enumerate}
\item for each $b\in B$ there is $a\in A$ such that
$b\iota=\iota a$;
\item a member of $B$ fixing $\iota(V)$ pointwise is the identity;
\item every $a\in A$ has an extension $b\in B$ satisfying
$b\iota=\iota a$;
\item for such a pair $(b,a)$, one has $b\in B_0$ if and only if $a\in A_0$.
\end{enumerate}
Then restriction induces an isomorphism $B\cong A$ carrying $B_0$ onto $A_0$.
Every extension in the third condition is unique, and the resulting extension
map is the inverse homomorphism.
\end{theorem}
\begin{proof}
The injection $\iota$ makes the element $a$ in the first condition unique;
write $\rho(b)=a$. For $b,c\in B$,
\[
 bc\iota=b\iota\rho(c)=\iota\rho(b)\rho(c),
\]
so uniqueness gives $\rho(bc)=\rho(b)\rho(c)$. The same argument gives
$\rho(1)=1$. Its kernel consists exactly of the elements fixing $\iota(V)$
pointwise, and is trivial by the second condition. The third condition makes
$\rho$ surjective. The last condition gives
$\rho(B_0)=A_0$: containment follows immediately, while a preimage of any
$a\in A_0$ belongs to $B_0$. Thus $\rho$ is the asserted pair isomorphism.
Injectivity gives uniqueness of extensions, and the inverse of an isomorphism
is a homomorphism.
\end{proof}

For colored graphs $X,Y$, one applies this theorem to $A=C(X)$,
$B=C(Y)$ and their palette kernels. Recognizing a retained copy of $V(X)$
is not enough: the restricted permutation must preserve the \emph{reconstructed}
old incidences and colors. In particular, old edges need not remain edges of
$Y$. Conversely, an explicit action of $C(X)$ on $Y$ establishes only the
extension condition, until rigidity rules out further correlations.

Here is a local criterion for that rigidity. All induced graphs and boundaries
in its statement are taken in the constructed graph $Y$.

\begin{lemma}[Recognition by component boundaries]\label{UR.Components}
Let $S\subseteq V(Y)$ be invariant under a group $B$ of correlations of $Y$.
For each connected component $D$ of $Y\setminus S$, define its boundary
\[
 \partial D=\{s\in S:\text{$s$ has a neighbor in $D$}\}.
\]
Assume distinct components have distinct boundaries. Suppose also that any
correlation of the induced colored graph on $D\cup\partial D$ fixing
$\partial D$ pointwise fixes $D$ pointwise. Then a member of $B$ fixing $S$
pointwise is the identity on $Y$.
\end{lemma}
\begin{proof}
Such a member permutes the components of $Y\setminus S$ and sends each
boundary to itself. Distinctness of the boundary sets therefore fixes each
component setwise. Its restriction to $D\cup\partial D$ is a correlation,
fixes the boundary pointwise, and is the identity by the local assumption.
This holds for every component; it is already the identity on $S$.
\end{proof}

Only color classes occurring in an induced subgraph are used when speaking of
its restricted coloring. A global correlation preserving that subgraph does
permute those classes, even if it permutes other colors outside the subgraph.
Thus the local hypothesis in Lemma~\ref{UR.Components} does not silently
replace correlations by color-preserving automorphisms.
